\section{General SPD self-adjointness}\label{sec:spd}
Because ``symmetrizable'' often denotes positive diagonal symmetrization in matrix literature, we use \emph{diagonal balance} for that stronger node-local property and \emph{SPD self-adjointness} for \(GJ=J^TG\) with arbitrary \(G>0\). See, for example, the positive-diagonal convention in \cite[Definition~4.1]{MS}.

\begin{theorem}[Finite real criterion]\label{thm:spd}
A real square J admits a real symmetric positive-definite G with \(GJ=J^TG\) if and only if J is diagonalizable over \(\R\) with entirely real spectrum.
\end{theorem}
\begin{proof}
If G exists, \(Q=G^{1/2}\) makes \(QJQ^{-1}\) real symmetric, since
\((QJQ^{-1})^T=Q^{-1}J^TQ=QJQ^{-1}\).
Conversely, write \(J=V\Lambda V^{-1}\) with V real invertible and \(\Lambda\) real diagonal. For any SPD K commuting with \(\Lambda\),
\begin{equation}
G=V^{-T}KV^{-1}\label{eq:SPDconstruct}
\end{equation}
is positive and balances J. These are all solutions: multiplying balance by \(V^T,V\) gives \(K\Lambda=\Lambda K\).
\end{proof}
Repeated semisimple roots and zero eigenvalues are allowed. An alternative, avoiding basis choices inside eigenspaces, is
\begin{equation}
G=\sum_\rho c_\rho E_\rho^TE_\rho,\qquad
E_\rho=\prod_{\sigma\ne\rho}\frac{J-\sigma I}{\rho-\sigma},\qquad c_\rho>0 .
\end{equation}
Here \(\sum E_\rho=I\) proves positivity, and \(E_\rho J=\rho E_\rho\) proves balance. In general these forms mix nodes and are nonunique. Positive diagonal detailed balance is strictly stronger than arbitrary SPD self-adjointness. Neither existence nor nonuniqueness selects a physical geometry. \src{GR2 \S2, L17; elementary lemma, no priority claim.}

\section{Exact positive-triad spectrum}\label{sec:triad}
Define
\begin{equation}
\sigma=a+b+c,\qquad
h=\sigma(1/a+1/b+1/c)-3,\qquad
t=\frac{\sigma(a^2+b^2+c^2)}{abc},\qquad d=1-3\ell .
\end{equation}
The invariant t is not time. The symmetric matrix
\begin{equation}
M_3=\sigma R^{-1}-\one\one^T,\qquad H_r=I-gM_3
\end{equation}
has kernel r. Weighted Cauchy--Schwarz gives
\[
v^TM_3v=\sigma\sum_i v_i^2/r_i-(\sum_i v_i)^2\ge0,
\]
with equality precisely for \(v\parallel r\). Its positive eigenvalues \(\mu_\pm\) therefore have sum h and product t, by trace and principal minors.

\begin{theorem}[Reality of the transverse quotient]\label{thm:triad}
The characteristic polynomial is
\begin{align}
\det(\xi I-J)&=(\xi-1)Q(\xi),\\
Q(\xi)&=\xi^2-d(2-gh)\xi+d^2(1-gh+g^2t),\label{eq:triadpoly}\\
\mathcal D_3&=\boxed{d^2g^2(h^2-4t)\ge0},\quad
\rho_\pm=d(1-g\mu_\pm).\label{eq:triadD}
\end{align}
The transverse quotient is always real semisimple. Its roots are distinct exactly when \(gd\ne0\) and the positive radii are not all equal.
\end{theorem}
\begin{proof}
The transverse object is \(\R^3/\operatorname{span}\{\one\}\); the Euclidean sum-zero plane need not be J-invariant. P is similar to symmetric \(I-gM_3\); its noncommon eigenvectors span this quotient. The synchronizer \(dI+\ell\one\one^T\) acts as scalar d on it, proving the polynomial and semisimplicity.

The two positive eigenvalues of \(M_3\) coincide only at equal radii. If both equal \(\mu\), then \(M_3=\mu(I-rr^T/(r^Tr))\). Its off-diagonal entries are all \(-1\), forcing ab=ac=bc and hence equality. The converse follows by substitution. Cases g=0 and d=0 yield scalar quotient action.
\end{proof}

\subsection{Full-matrix defects and local SPD geometry}
Let \(w_\alpha\perp r\) be eigenvectors of \(M_3\), \(v_\alpha=R^{-1}w_\alpha\), and \(p_\alpha=1-g\mu_\alpha\). In the basis \((\one,v_-,v_+)\),
\begin{equation}
J\sim
\begin{pmatrix}
1&\ell p_-\one^Tv_-&\ell p_+\one^Tv_+\\
0&dp_-&0\\0&0&dp_+
\end{pmatrix}.\label{eq:triadres}
\end{equation}
This follows by applying P and then \(dI+\ell\one\one^T\) to each vector. Thus, for nonconstant radii and \(g\ell d\ne0\), a defect occurs exactly when
\begin{equation}
dp_\alpha=1,\qquad \one^Tv_\alpha\ne0.\label{eq:resonance}
\end{equation}
It is a size-two Jordan block at one, not a defective transverse coincidence.

To classify the vanishing coupling, multiply \(M_3w=\mu w\) by \(\one^T\):
\[
\one^TR^{-1}w=(\mu+3)\one^Tw/\sigma .
\]
If this is zero, the eigen-equation reduces to \(\sigma R^{-1}w=\mu w\). A nonzero sum-zero w can then be supported only on equal radii. Consequently every resonance for three distinct radii is defective; for exactly two equal radii the antisymmetric equal-pair mode is semisimple at resonance and the other mode is defective. Equal radii and cases \(g=0,\ell=0,d=0\) are diagonalizable; at d=0, \(J^2=J\).

Near a strictly stable equal background the transverse eigenvalues remain inside the unit disk and away from one. These resonances are excluded, so Theorem~\ref{thm:spd} supplies an SPD metric throughout a sufficiently small positive unequal triad branch, even when Theorem~\ref{thm:cycle} excludes every positive diagonal weight.

Finally, let \(v_i=u_i/\sqrt{k}\) in the weak expansion. The first variation of \(M_3\), restricted to the equal-background transverse plane, is
\(-3\Pi\diag(v)\Pi\), \(\Pi=I-\one\one^T/3\). It is traceless and has squared eigenvalue gap \(6\sum v_i^2\). Therefore
\begin{equation}
\mathcal D_3=
6d^2g^2\left(\frac{\eps}{2\eps k+3g}\right)^2
\sum_i\kappa_i^2\,\eta^2+O(\eta^3).\label{eq:triadweak}
\end{equation}
The positive quadratic onset splits real roots; no order makes the positive-triad transverse spectrum complex. This is independent of the cubic cycle obstruction. \src{GR2 \S\S3--4, L18--L20.}

\section{Exact finite Bloch decomposition}\label{sec:bloch}
Number nodes by cell \(n=0,\ldots,q-1\) and residue \(\alpha=0,1,2\). For
\(\phi_{n,\alpha}=z^nv_\alpha\), \(z=e^{i\theta}\), \(\theta=2\pi m/q\),
\begin{align}
D(z)&=\begin{pmatrix}-2&1&z^{-1}\\1&-2&1\\z&1&-2\end{pmatrix},\\
H_r(z)&=I-g\diag((\sigma-3r_\alpha)/r_\alpha)+gD(z),\\
J(z)&=\boxed{(I+\ell D(z))R^{-1}H_r(z)R}.\label{eq:bloch}
\end{align}
Substitution at cell boundaries supplies the \(z^{\pm1}\) factors. Fourier transformation in n is an exact finite direct sum. Both D and \(H_r\) are Hermitian on the unit circle, but the product in \eqref{eq:bloch} need not be. Conjugate phases give conjugate blocks.

To prevent a collision of notation with the chiral scalar below, write
\begin{align}
V_r&=(a-b)(a-c)(b-c),\quad E_B=\ell(1+3g)-g,\\
\Gamma_B&=\ell gE_B V_r/(abc),\qquad w=z+z^{-1}-2,\\
C_B&=\ell g\left[\ell gt-\tfrac12(g+\ell-\ell g)h\right].
\end{align}
\begin{proposition}[Exact coefficients]\label{prop:blochpoly}
Writing \(\det(\xi I-J(z))=\xi^3-T_B\xi^2+B_B(z)\xi-D_B(z)\),
\begin{align}
T_B&=3-6\ell-dgh,\label{eq:BT}\\
B_B(z)&=d(2-gh)+d^2(1-gh+g^2t)+C_Bw-\frac{\Gamma_B}{2}(z-z^{-1}),\label{eq:BB}\\
D_B(z)&=(d^2+\ell^3w)(1-gh+g^2t+g^3w).\label{eq:BD}
\end{align}
\end{proposition}
\begin{proof}
Take the trace of \eqref{eq:bloch} to obtain \(T_B\). The middle coefficient is \(((\tr J)^2-\tr J^2)/2\). Collecting \(z,z^{-1}\) gives mean coefficient \(C_B\), difference \(-\Gamma_B\); its value at z=1 is fixed by \eqref{eq:triadpoly}. For the determinant, multiply
\(\det(I+\ell D)=d^2+\ell^3w\) and
\(\det H_r=1-gh+g^2t+g^3w\).
Similarity by R changes neither determinant. These operations derive all coefficients; Appendix~\ref{app:bloch} records a reproducible algebra route.
\end{proof}

\begin{theorem}[Blockwise spectral decision]\label{thm:blocks}
If \(\Gamma_B\sin\theta\ne0\), the block has non-real spectrum and the full real ring has no SPD symmetrizer. If this product vanishes, its coefficients are real and the cubic discriminant
\begin{equation}
\mathfrak D=T_B^2B_B^2-4B_B^3-4T_B^3D_B-27D_B^2+18T_BB_BD_B\label{eq:cubicD}
\end{equation}
classifies simple roots: positive means three distinct real roots; negative means a real root and a conjugate pair. At zero discriminant, a Jordan test is additionally necessary.
\end{theorem}
\begin{proof}
\(\im B_B=-\Gamma_B\sin\theta\) whereas \(T_B,D_B\) are real. An all-real root set would have real coefficients, a contradiction. In fact a real root would obey \((\im B_B)\xi=0\), so if \(D_B\ne0\) all three block roots are non-real. The conjugate block supplies their conjugates.

For a real cubic the discriminant is the product of squared root differences. It is positive for three distinct real roots; for a conjugate pair the squared difference of that pair is negative and the remaining paired product positive, giving negative sign. Repeated roots give zero. For a double root \(\rho\) and distinct \(\sigma_0\), diagonalizability is equivalent to \(\rank(J-\rho I)=1\), or \((J-\rho I)(J-\sigma_0I)=0\). At a triple root it is equivalent to J being scalar. This follows directly from the eigenspace dimensions or minimal polynomial.
\end{proof}
When not all roots coincide, the repeated root at zero discriminant is
\[
\rho=\frac{T_BB_B-9D_B}{2(T_B^2-3B_B)},\qquad \sigma_0=T_B-2\rho .
\]
Triple coincidence occurs when \(B_B=T_B^2/3,D_B=T_B^3/27\). If coefficients are non-real, repeated-root rank tests still distinguish complex Jordan defects, although SPD balance is already excluded.

The full ring admits an SPD symmetrizer exactly when every block has real semisimple spectrum. Conjugate Fourier eigenvectors then give a real eigenbasis; repeated roots in separate blocks do not create Jordan coupling.

For q at least three, an allowed phase has nonzero sine. Thus three distinct radii and
\begin{equation}
g\ell[\ell(1+3g)-g]\ne0\label{eq:generic}
\end{equation}
force non-real spectrum on every admitted \(N\ge9\), however weak the nonzero inequality. The immediate obstruction vanishes when g=0, \(\ell=0\), \(E_B=0\), \(V_r=0\), or the particular phase is \(z=\pm1\). Its vanishing is not sufficient for real semisimple spectrum.

Along the weak branch,
\[
\Gamma_B=\ell gE_B\left(\frac{\eps}{2\eps k+3g}\right)^3V(\kappa)\eta^3+O(\eta^4).
\]
At a simple nonzero equal-background block root rho, differentiating the characteristic equation gives
\[
\im\delta\rho=\frac{\Gamma_B\sin\theta\,\rho}
{\prod_{\sigma'\ne\rho}(\rho-\sigma')}+O(\eta^4).
\]
Repeated or zero baseline roots require the exact polynomial tests instead. Complex roots can remain strictly inside the unit disk. They are oscillatory linear-response modes, with no physical wave interpretation. \src{GR2 \S5, L21--L24.}

\section{Six-ring reality, defect and complex unfolding}\label{sec:six}
For N=6 only \(z=1,-1\) occur; N=12 adds \(z=i,-i\). The first block is the triad. At z=-1 the polynomial is real with w=-4; at \(z=\pm i\), its imaginary middle coefficient is \(\mp\Gamma_B\). Thus the same native background can have a real-diagonalizable six-ring response and a non-real twelve-ring response.

At equal radii the antiperiodic block has double root
\(\rho_1=(1-\ell)(1-g)\) and simple root
\(\rho_4=(1-4\ell)(1-4g)\), with gap \(3E_\pi\), \(E_\pi=g+\ell-5g\ell\).
For \(v_i=u_i/\sqrt{k}\), its discriminant is
\begin{equation}
\mathfrak D_\pi=
6g^2(1-\ell)^2\sum_i v_i^2(3E_\pi)^4\eta^2+O(\eta^3).\label{eq:sixgeneric}
\end{equation}
Indeed the first derivative is
\((I+\ell D_\pi)g([D_\pi,\diag v]+3\diag v)\).
Its projection onto the double \(-1\) eigenspace of \(D_\pi\) is
\(3g(1-\ell)\Pi_\pi\diag(v)\Pi_\pi\), symmetric with the squared gap shown in \eqref{eq:sixgeneric}. Multiply by the fourth power of separation from the third root. At a fixed stable baseline with \(g(1-\ell)E_\pi\ne0\), weak nontrivial inequality therefore gives real simple antiperiodic roots and an SPD full-ring metric, despite the diagonal obstruction.

The neighbourhood is not uniform near exceptional baselines. At \(E_\pi=0\), \(\ell=g/(5g-1)\), the first variation \(K_\pi\) has trace and determinant zero and
\[
\tr K_\pi^2=\frac{3g^2(4g-1)^2}{(5g-1)^2}\sum_i v_i^2.
\]
Unless this coefficient vanishes, its three first-order roots are distinct and real. At \(\ell=1\), the antiperiodic synchronizer has rank one. The exact response can instead be defective.

\begin{theorem}[Exact native nilpotent response and unfolding]\label{thm:nilpotent}
Take
\begin{equation}
r=(1-\tau,1,1+\tau),\quad g=(1-\tau^2)/4,\quad
\ell=1,\quad \eps=1/20,\quad
k_i=r_i^2-\frac g\eps\frac{(\Delta_3r)_i}{r_i},\label{eq:fixture}
\end{equation}
with \(0<|\tau|\) sufficiently small. Then
\begin{equation}
\boxed{J(-1)=\frac{\tau^2}{4}
\begin{pmatrix}1\\-1\\1\end{pmatrix}
\begin{pmatrix}1+\tau&2&1-\tau\end{pmatrix}.}\label{eq:nilpotent}
\end{equation}
It has rank one, square zero, and Jordan sizes \(2+1\) at eigenvalue zero. Holding r,g,\(k_i,\eps\) fixed and varying only \(\ell=1+\delta\), its discriminant satisfies
\begin{equation}
\mathfrak D_\pi=\frac{27}{2}\tau^{12}\delta^3+O(\delta^4).\label{eq:unfold}
\end{equation}
For each fixed nonzero small tau, sufficiently small negative delta gives a complex pair, positive delta gives three distinct real roots, and zero delta gives the defect.
\end{theorem}
\begin{proof}
Substitute \eqref{eq:fixture} in \eqref{eq:bloch}; direct multiplication gives \eqref{eq:nilpotent}. Its nonzero row and column pair to zero, proving rank one and square zero. A three-dimensional nonzero square-zero rank-one operator has one size-two block and one size-one block. The coefficients under the stated variation are exactly
\begin{align}
T_B&=-3\delta/2,\\
B_B&=-\frac{3\delta}{16}\,[8\tau^4+\delta(10\tau^4+\tau^2-3)],\\
D_B&=\frac{\delta^2\tau^4(4\delta+3)(9-\tau^2)}{16}.\label{eq:unfoldcoeff}
\end{align}
Substituting into \eqref{eq:cubicD} leaves the cubic coefficient \(27\tau^{12}/2>0\), proving the sign conclusions and simplicity away from zero for small delta.
\end{proof}

This is a native fixed point by construction, not an arbitrary matrix. The coefficients are positive for small tau, with mean \(k=1-28\tau^2/3\). For each fixture, \(\eta=\tau\) and \(\kappa_i=(k_i-k)/\tau\) put it in the admitted form. Varying tau in this auxiliary family does \emph{not} keep kappa fixed. At the rational fixture \(\tau=10^{-3}\), exact contraction-box inequalities connect the point to the unique weak branch for its fixed direction (Appendix~\ref{app:fixture}).

The equal-reference noncommon phase multipliers at \(\ell=1\) are \(0,-1/2-3\tau^2/2,-3\tau^2\); its radial multipliers are strictly inside the unit disk for small tau. At the unequal point the triad transverse roots are
\(-\tfrac12(1\pm\tau\sqrt{3+6\tau^2})\), and the antiperiodic eigenvalues vanish. Moreover \(\|J(-1)\|_\infty=\tau^2\). Thus the defect and its sufficiently small unfolding can remain normally stable. Stability, reality and diagonalizability are separate questions. \src{GR2 \S6, L25--L27.}

\begin{figure}[htbp]\centering
\includegraphics[width=\linewidth]{../figures/03_spectral_hierarchy.pdf}
\caption{Finite spectral hierarchy, with hypotheses stated on each branch. The six-ring three-way alternative uses \eqref{eq:fixture} and variation of the existing synchronizer parameter only. The larger-ring statement requires \eqref{eq:generic} and distinct radii. All are finite linear-response results. Source: GR2 \S\S3--6.}
\end{figure}
